\documentclass[11pt]{article}

\usepackage[margin=1in]{geometry}
\usepackage{amsmath}
\usepackage{amssymb}
\usepackage{array}
\usepackage{booktabs}
\usepackage{hyperref}
\usepackage{listings}
\usepackage{xcolor}

\hypersetup{
  colorlinks=true,
  linkcolor=blue,
  citecolor=blue,
  urlcolor=blue
}

\lstset{
  basicstyle=\ttfamily\small,
  breaklines=true,
  frame=single,
  columns=fullflexible
}

\title{Accelerating Gradient Boosting Decision Trees via Hardware Ray Tracing Cores and On-Device Fused Reduction}
\author{GAFIME Project Technical Disclosure Draft}
\date{\today}

\begin{document}

\maketitle

\begin{abstract}
Gradient boosting decision trees (GBDTs) and tree-derived decision-path features
are central to tabular machine learning, but scoring large sets of path
candidates remains dominated by candidate-row predicate evaluation, intermediate
membership materialization, and reduction of the resulting binary indicators
against a target. This report describes a hardware mapping for shallow
axis-aligned decision-path candidates onto NVIDIA RTX hardware ray tracing
cores through OptiX. A decision path such as $x_j > a \land x_k \le b$ is encoded
as a bounded axis-aligned region, rows are encoded as points, and candidate
membership becomes point-in-region traversal. The implementation in GAFIME v1
keeps Python and Rust responsible for public API, validation, planning,
scheduling, and result ownership, while CUDA owns the RT-specific execution
boundary: finite box planning, triangle/AABB geometry construction, OptiX
pipeline launch, exact any-hit predicate guards, device-side compact scoring,
and explicit unsupported behavior when the RT contract cannot be proven.

The key performance path is not path-major membership export. Instead, the CUDA
payload evaluates compact decision-path scores through optional
\texttt{gafime\_gpu\_decision\_path\_score}: OptiX traversal feeds on-device
counts and target sums, and a fused reduction writes compact result rows of size
\texttt{paths} $\times$ \texttt{metrics}. For tree-leaf-like batches whose
regions are finite, bounded, two-dimensional, and non-overlapping, a
\texttt{firsthit} mode terminates each ray after the first exact in-box hit and
fails closed when the non-overlap proof is unavailable.

On an NVIDIA GeForce RTX 4060 Laptop GPU (Ada, sm\_89), the parity-covered
partitioned first-hit benchmark evaluates $262{,}144 \times 8{,}192 =
2.147$ billion membership-equivalent candidate-row tests in $0.88$--$0.94$ ms,
or approximately $2.28$--$2.45$ trillion evaluations/s, with maximum absolute
score difference $1.29454 \times 10^{-7}$ against a CPU reference. A larger
throughput-only run evaluates $262{,}144 \times 1{,}048{,}576 = 274.878$
billion membership-equivalent tests in $20.030$ ms, or $13.724$ trillion
evaluations/s; this larger run is performance evidence and is paired with
smaller parity-covered runs. The result suggests that fixed-function traversal
hardware can accelerate a specific, practically important subset of GBDT-like
decision-path scoring when candidate regions form non-overlapping spatial
partitions and output remains compact.
\end{abstract}

\section{Disclosure Status}

This document is a technical disclosure draft for the GAFIME v1 CUDA RT-core
decision-path implementation. It is written in an arXiv-compatible article
format and is intended to preserve the architecture, assumptions, numerical
contract, and measured evidence behind the work. It is not a legal opinion, a
patent filing, or a claim that the full research validation expected of a
conference submission has been completed. The implementation and measurements
described here correspond to the GAFIME branch
\texttt{codex/rt-gbdt-cuda}, through commit
\texttt{5fe114a test(cuda): add first-hit RT scale proof}.

\section{Introduction}

GBDTs remain among the strongest practical methods for tabular data. The
original gradient boosting view of Friedman \cite{friedman2001gbm} treats
boosting as function-space optimization, while systems such as XGBoost
\cite{chen2016xgboost}, LightGBM \cite{ke2017lightgbm}, and CatBoost
\cite{prokhorenkova2018catboost} demonstrated that implementation details such
as cache layout, sparsity handling, histogram construction, categorical feature
handling, and sampling strategies are central to practical performance.

This work focuses on a different but related bottleneck: scoring many
tree-derived decision-path candidates as explicit features or interaction
candidates. A path from a tree root to an internal node or leaf is a conjunction
of threshold predicates:
\[
  (x_{f_1} \le t_1) \land (x_{f_2} > t_2) \land \cdots.
\]
For low-depth paths over one to three numerical features, the conjunction is
equivalent to membership in an axis-aligned interval, rectangle, or box. In a
candidate-mining engine, evaluating such paths over many rows and many candidate
regions creates a large candidate-row workload. Naively materializing a
\texttt{rows} $\times$ \texttt{paths} membership matrix is memory-expensive, and
branch-heavy scalar or SIMT predicate loops can underutilize the available
device.

Hardware ray tracing cores are designed to answer geometric membership and
intersection queries at high throughput. NVIDIA OptiX is a programmable ray
tracing system for GPUs \cite{parker2010optix}, and RTX/Turing-class GPUs
introduced dedicated RT cores for ray traversal and intersection acceleration
\cite{nvidia2018turing}. Prior ray traversal literature studied the mapping of
traversal and intersection to GPU execution \cite{aila2009raytraversal,
aila2012raytraversal}, while ray tracing kernel frameworks such as Embree
\cite{wald2014embree} demonstrated that specialized traversal kernels can be
valuable systems primitives beyond image synthesis.

The central observation of this report is simple:

\begin{quote}
  A shallow GBDT decision path is an axis-aligned geometric region, and a row is
  a point. Therefore, scoring many path candidates can be cast as a batched
  point-in-region traversal problem.
\end{quote}

The practical contribution is not merely this geometric analogy. The GAFIME
implementation makes the mapping operational under a strict systems contract:
Rust owns orchestration and candidate semantics; CUDA owns RT execution; the
public result remains compact; unsupported RT cases fail explicitly; and the
measured fast path is tied to parity-covered smaller runs.

\subsection{Contributions}

This draft documents the following contributions.

\begin{enumerate}
  \item A mapping from finite, bounded, shallow decision-path predicates to
  OptiX triangle/AABB geometry while preserving exact GAFIME predicate
  semantics through any-hit guards.
  \item A CUDA C ABI scoring path,
  \texttt{gafime\_gpu\_decision\_path\_score}, that avoids host materialization
  of path-major membership and emits compact \texttt{ResultTable} rows.
  \item A grouped, instanced OptiX launch for many feature-pair groups: each
  group builds its own triangle GAS, the launcher wraps them in one IAS, and
  ray generation launches \texttt{rows} $\times$ \texttt{groups} rays.
  \item A fail-closed first-hit mode for tree-leaf-like non-overlapping
  partitions, using \texttt{OPTIX\_RAY\_FLAG\_TERMINATE\_ON\_FIRST\_HIT} and an
  explicit non-overlap proof before allowing the fast path.
  \item A release-measure proof script that enforces both minimum throughput and
  maximum numerical deviation for the parity-covered first-hit workload.
\end{enumerate}

\section{Problem Formulation}

Let $X \in \mathbb{R}^{n \times d}$ be a tabular matrix with $n$ rows and $d$
features, and let $y \in \mathbb{R}^{n}$ be a target vector. A decision-path
candidate $p$ is represented as a sequence of terms
\[
  p = \{(f_\ell, s_\ell, t_\ell)\}_{\ell=1}^{L},
\]
where $f_\ell$ is a feature index, $s_\ell \in \{\le, >\}$ is a split sign, and
$t_\ell$ is a threshold. The binary membership value for row $i$ is
\[
  z_{p,i} =
  \begin{cases}
    1 & \text{if all terms in } p \text{ hold for row } i, \\
    0 & \text{otherwise}.
  \end{cases}
\]
The compact scoring path currently supports Pearson correlation and $R^2$.
For a path $p$, define:
\[
  S_x = \sum_i z_{p,i}, \quad
  S_y = \sum_i y_i, \quad
  S_{yy} = \sum_i y_i^2, \quad
  S_{xy} = \sum_i z_{p,i} y_i.
\]
The centered covariance terms are:
\[
  C_{xx} = S_x - \frac{S_x^2}{n}, \quad
  C_{yy} = S_{yy} - \frac{S_y^2}{n}, \quad
  C_{xy} = S_{xy} - \frac{S_x S_y}{n}.
\]
The Pearson score is:
\[
  \rho_p =
  \begin{cases}
    \mathrm{clip}\left(\frac{C_{xy}}{\sqrt{C_{xx} C_{yy}}}, -1, 1\right)
      & \text{if } C_{xx} C_{yy} > 0, \\
    0 & \text{otherwise}.
  \end{cases}
\]
The $R^2$ score is $\rho_p^2$ clamped to $[0,1]$.

The computational objective is to score $m$ paths over $n$ rows without
materializing an $n \times m$ floating membership matrix. We refer to $n \cdot m$
as the number of membership-equivalent evaluations, or candidate-row
evaluations. This unit is used for throughput because it measures the logical
amount of predicate membership work implied by the score computation.

\section{Geometric Mapping}

Each decision-path term tightens a lower or upper bound along a feature axis.
For a path over at most three unique features, the term set maps to an
axis-aligned box:
\[
  B_p = [l_1, h_1] \times [l_2, h_2] \times [l_3, h_3],
\]
with lower-open bounds for strict $>$ terms and upper-closed bounds for $\le$
terms. A row becomes a point
\[
  q_i = (X_{i,f_1}, X_{i,f_2}, X_{i,f_3}).
\]
Membership is therefore equivalent to $q_i \in B_p$ under the open/closed
predicate convention.

The implementation uses two geometry modes:

\begin{itemize}
  \item \textbf{Bounded 2D triangle geometry.} A finite two-dimensional box is
  represented as two triangles. This path exercises the fixed-function triangle
  traversal hardware. The any-hit program still rechecks the exact GAFIME
  predicate semantics, including lower-open and upper-closed boundaries.
  \item \textbf{Custom AABB geometry.} One-dimensional, three-dimensional, or
  open-bound cases use exact custom AABB intersection when representable.
\end{itemize}

The first-hit fast path applies only to the bounded 2D case and only when every
internal group is proven non-overlapping. This condition matches tree-leaf
partitions: within one tree or one feature-pair group, a row belongs to at most
one leaf region.

\section{System Architecture}

\subsection{Ownership Boundary}

GAFIME v1 follows a hard ownership split. Python exposes the public API, Rust
owns validation, candidate planning, backend selection, scheduling, and compact
report ownership, while CUDA owns CUDA-specific execution. The RT-core
implementation respects this contract:

\begin{itemize}
  \item Rust passes compact \texttt{GafimeDecisionPathTerm} and path-offset
  buffers through the GPU C ABI.
  \item CUDA validates RT representability, builds geometry, launches OptiX, and
  writes compact metric rows.
  \item CUDA does not discover paths, change scheduler policy, infer Python
  behavior, or silently fall back across backends.
  \item Unsupported RT requests return explicit unsupported status when the
  caller requires RT or when first-hit mode is requested.
\end{itemize}

\subsection{CUDA Source Layout}

The CUDA RT implementation is explicitly separated from the generic CUDA
continuous kernels:

\begin{itemize}
  \item \texttt{src/cuda/rt\_kernels.cu}: OptiX device programs, RT-specific
  CUDA kernels, point packing, membership kernels, and direct-score reductions.
  \item \texttt{src/cuda/rt\_launcher.cu}: host-side RT planning, finite-box
  validation, OptiX pipeline setup, GAS/IAS construction, cached workspaces,
  grouped execution, explicit SM comparator dispatch, and score dispatch.
  \item \texttt{src/cuda/launcher.cu}: public C ABI bridge only for decision
  path RT symbols.
\end{itemize}

This matters for maintainability: RT-specific traversal logic is not hidden
inside the generic CUDA continuous scoring path.

\subsection{Compact Score ABI}

The public compact score symbol is:

\begin{lstlisting}
gafime_gpu_decision_path_score(matrix, paths, result_out)
\end{lstlisting}

The input \texttt{paths} buffer contains:

\begin{itemize}
  \item path count,
  \item term count,
  \item term array,
  \item path offsets,
  \item metric ids,
  \item metric count,
  \item flags such as \texttt{GAFIME\_DECISION\_PATH\_FLAG\_REQUIRE\_RT}.
\end{itemize}

The output is a compact \texttt{GafimeResultTable}. It has one row per path,
not one row per path-row membership value. This is the main memory difference
from path-major membership materialization.

\subsection{Grouped Instanced OptiX Launch}

Mixed-axis workloads may contain paths over feature pairs $(f_0,f_1)$,
$(f_2,f_3)$, $(f_1,f_2)$, and so on. A single OptiX GAS cannot directly encode
all possible feature-axis coordinate systems without changing the point
interpretation. GAFIME groups paths by compatible axes, builds one triangle GAS
per group, and wraps the GAS handles in an instance acceleration structure
(IAS). Ray generation then launches:
\[
  \texttt{launch dimensions} = (n, g, 1),
\]
where $g$ is the number of groups. Group-local points are prepacked as $x,y$
pairs. The any-hit program uses the OptiX instance id plus a group path-offset
table to recover the flattened path id, and a final scatter map restores public
result order.

This design gives OptiX a large launch while keeping Rust as the owner of
candidate discovery and global scheduling.

\subsection{Scoring Modes}

The implementation has three score paths:

\paragraph{Bitset score path.}
OptiX writes an idempotent device bitset. A CUDA kernel reduces the bitset with
the target vector. This has tighter deterministic parity than direct float
atomics and is the safer default path.

\paragraph{Direct score path.}
OptiX any-hit updates per-path counts and target sums directly with device
atomics. This removes the temporary bitset and reduction scan, but floating
atomic accumulation is not bit-stable. It remains opt-in with documented
tolerance.

\paragraph{First-hit direct path.}
For non-overlapping finite bounded 2D groups, any-hit records the first exact
in-box hit and terminates the ray. This avoids duplicate triangle ambiguity and
models tree-leaf partition semantics. It is selected with:

\begin{lstlisting}
GAFIME_CUDA_DECISION_PATH_RT_SCORE=firsthit
\end{lstlisting}

First-hit mode is fail-closed: if CUDA cannot prove the batch is finite,
bounded, two-dimensional, and non-overlapping inside every RT group, the score
ABI returns unsupported instead of silently routing the opt-in request to the SM
comparator.

\subsection{Caching and Residency}

The CUDA RT launcher caches:

\begin{itemize}
  \item host grouped plans,
  \item group GASes and IAS,
  \item OptiX pipelines and workspaces,
  \item prepacked grouped points keyed by feature pointer, generation, row
  count, and geometry signature,
  \item target-wide statistics keyed by target pointer, generation, and row
  count,
  \item flattened original-path scatter maps,
  \item persistent grouped scratch buffers.
\end{itemize}

Target-only updates invalidate target statistics but not feature-derived packed
points. Feature uploads invalidate feature-derived caches. This makes repeated
resident-session scoring substantially different from rebuilding the whole
pipeline per public call.

\section{Numerical Contract}

The correctness policy separates logical membership from floating score
accumulation.

\begin{itemize}
  \item Membership semantics are exact: \texttt{LE} means $x \le t$,
  \texttt{GT} means $x > t$.
  \item NaN handling is not silently changed. Uploaded finite feature matrices
  are required for the RT path; otherwise RT returns unsupported.
  \item First-hit mode is permitted only for non-overlapping 2D groups. This is
  the condition that makes first-hit membership equivalent to leaf assignment.
  \item Integer counts, path ids, and compact row ordering must be exact.
  \item Floating Pearson/R2 values may differ by documented tolerance when
  direct float atomics or backend instruction differences make bit parity
  impractical.
\end{itemize}

The current parity-covered first-hit run has maximum absolute difference
$1.29454 \times 10^{-7}$ against the CPU reference, far below the documented
spike tolerance of $10^{-4}$.

\section{Experimental Setup}

The primary local system used for the measurements documented in this draft is:

\begin{itemize}
  \item CPU: AMD Ryzen AI 9 HX 370 class laptop CPU, with AVX512 path compiled
  into the benchmark reference.
  \item GPU: NVIDIA GeForce RTX 4060 Laptop GPU, Ada generation, sm\_89,
  approximately 7.63 GiB global memory reported by the benchmark.
  \item CUDA payload: GAFIME v1 CUDA C ABI shared library built with OptiX RT
  enabled.
  \item Benchmark: \texttt{tests/gpu/cuda\_rt\_membership\_scale\_bench.cpp}.
  \item Release-measure wrapper:
  \texttt{tests/release\_measure/perf\_05\_cuda\_rt\_firsthit\_scale.py}.
\end{itemize}

The benchmark uses synthetic bounded 2D boxes over feature-pair groups. The
partitioned-grid mode creates non-overlapping boxes within each feature-pair
group, matching the leaf-partition invariant needed by first-hit traversal.

\section{Results}

\subsection{Membership Materialization is Not the Main Win}

Path-major membership materialization remains expensive because it writes
\texttt{rows} $\times$ \texttt{paths} output. On a $1{,}048{,}576 \times 512$
workload, the measured RT membership path was in the same low-G evaluations/s
band as the SM comparator:

\begin{center}
\begin{tabular}{lrrrr}
\toprule
Mode & Rows & Paths & Time (ms) & G eval/s \\
\midrule
CPU AVX512 membership & 1,048,576 & 512 & 111.218 & 4.827 \\
CUDA RT membership & 1,048,576 & 512 & 194.236 & 2.764 \\
CUDA SM membership & 1,048,576 & 512 & 189.348 & 2.835 \\
\bottomrule
\end{tabular}
\end{center}

This result motivates compact scoring. The RT core should not merely create a
larger membership matrix faster; it should avoid materializing that matrix.

\subsection{Compact Scoring Improves the Workload Shape}

The compact score ABI reduces output to \texttt{paths} $\times$ \texttt{metrics}.
On a $1{,}048{,}576 \times 512$ workload:

\begin{center}
\begin{tabular}{lrrrr}
\toprule
Mode & Rows & Paths & Time (ms) & G eval/s \\
\midrule
CUDA RT compact score & 1,048,576 & 512 & 9.118 & 58.881 \\
CUDA SM compact score & 1,048,576 & 512 & 26.768 & 20.056 \\
\bottomrule
\end{tabular}
\end{center}

The temporary score mask for this workload is 64 MiB as a bitset, compared with
512 MiB as a byte mask and 2.00 GiB as an \texttt{f32} membership matrix.

\subsection{Direct Traversal Statistics}

Direct traversal statistics remove the bitset and reduction pass, but use
floating atomic accumulation. Representative direct-mode runs:

\begin{center}
\begin{tabular}{lrrrrr}
\toprule
Rows & Paths & Time (ms) & G eval/s & RT max abs diff & Notes \\
\midrule
1,048,576 & 512 & 3.195 & 168.048 & $5.04\times10^{-6}$ & direct \\
262,144 & 512 & 1.572 & 85.397 & $7.58\times10^{-6}$ & direct \\
262,144 & 8,192 & 12.445 & 172.555 & $4.10\times10^{-6}$ & grouped \\
\bottomrule
\end{tabular}
\end{center}

These numbers show a strong direction but are not the default deterministic
path because of floating atomic ordering.

\subsection{First-Hit Partitioned Path}

For non-overlapping tree-leaf-like partitions, first-hit direct scoring is the
main result.

\begin{center}
\begin{tabular}{lrrrrr}
\toprule
Case & Rows & Paths & Time (ms) & G eval/s & Parity \\
\midrule
Parity-covered first-hit & 262,144 & 8,192 & 0.877--0.942 & 2,278--2,450 &
$1.29454\times10^{-7}$ \\
Throughput-only first-hit & 65,536 & 1,048,576 & 18.059 & 3,805 & skipped \\
Throughput-only first-hit & 262,144 & 1,048,576 & 20.030 & 13,724 & skipped \\
\bottomrule
\end{tabular}
\end{center}

The most recent release-measure run on the current checkout produced:

\begin{lstlisting}
case: rows=262144 paths=8192 evals=2147.484M
resident upload        3.832 ms
cpu_score_ref      6170.278 ms     0.348 G eval/s
gpu_rt_score          0.942 ms  2278.804 G eval/s
score parity      rt_max_abs=1.29454e-07 sm_max_abs=0
\end{lstlisting}

The largest documented throughput-only run produced:

\begin{lstlisting}
rows=262,144 paths=1,048,576 partitioned-grid overlap-axis axis_pairs=8
gpu_rt_score firsthit 20.030 ms  13723.633 G eval/s
score parity          skipped (--throughput-only)
\end{lstlisting}

The throughput-only result should be interpreted as performance evidence for
the same implementation and workload family, not as independent correctness
evidence. Correctness is established by smaller parity-covered runs.

\section{Why the Result is Large}

The first-hit result is large because the workload aligns with the RT hardware
contract:

\begin{enumerate}
  \item The candidate regions are spatial boxes, not arbitrary predicates.
  \item The boxes are non-overlapping inside each group, so each row can
  terminate after the first exact in-box hit.
  \item The launch shape is large: \texttt{rows} $\times$ \texttt{groups} rays,
  rather than many tiny launches.
  \item The public result is compact: the output is metric rows, not a dense
  membership matrix.
  \item Repeated resident calls reuse geometry, packed points, target stats, and
  scratch buffers.
\end{enumerate}

Dense overlapping boxes do not show the same acceleration. In million-path
overlap stress cases, throughput was around 95--99 G eval/s and parity was
skipped at that size. This is consistent with the interpretation that RT cores
are most effective here when the candidate set resembles a tree partition rather
than a dense random overlap field.

\section{Limitations}

\begin{itemize}
  \item The fastest first-hit path currently applies only to finite, bounded,
  non-overlapping 2D groups.
  \item Compact RT decision-path score currently supports Pearson and $R^2$.
  MI and Spearman require additional parity work.
  \item CUDA/OptiX is the only RT-core backend in this spike. ROCm and Metal do
  not expose this RT path through the current GAFIME ABI.
  \item The 13.7 T eval/s result is throughput-only. It is paired with
  parity-covered smaller runs, but it is not itself CPU-parity validated.
  \item Measurements are from one laptop Ada GPU. Larger RTX devices may shift
  the bottleneck toward acceleration-structure size, atomic statistic pressure,
  launch scheduling, memory bandwidth, or RT-core occupancy.
  \item NCU did not expose a direct RT-core saturation percentage for the OptiX
  launch. Available counters showed that the visible limiter was not ordinary
  branch divergence, but further hardware-counter work is needed.
  \item This is decision-path scoring acceleration, not a complete replacement
  for GBDT training algorithms such as XGBoost, LightGBM, or CatBoost.
\end{itemize}

\section{Related Work}

Friedman's gradient boosting machine \cite{friedman2001gbm} established the
function-space optimization view of boosting. XGBoost \cite{chen2016xgboost}
emphasized scalable systems design for tree boosting, including cache access
patterns, compressed columns, sparsity-aware split finding, and sharding.
LightGBM \cite{ke2017lightgbm} introduced gradient-based one-side sampling and
exclusive feature bundling to improve GBDT training efficiency. CatBoost
\cite{prokhorenkova2018catboost} addressed ordered boosting and categorical
feature handling. These systems target training and inference efficiency through
histograms, sampling, layout, and algorithmic changes.

OptiX \cite{parker2010optix} provides programmable ray generation,
intersection, any-hit, closest-hit, and traversal composition over GPU
acceleration structures. RTX/Turing hardware introduced dedicated RT cores
\cite{nvidia2018turing}, making fixed-function traversal available to
applications through APIs such as OptiX and DXR. Ray traversal efficiency work
\cite{aila2009raytraversal,aila2012raytraversal} studied the practical
relationship between traversal, memory behavior, SIMD/SIMT utilization, and
scene structure. Embree \cite{wald2014embree} demonstrated the value of
optimized traversal kernels as reusable infrastructure on CPUs.

The present work differs from conventional ray tracing applications by using RT
traversal as a tabular candidate-scoring primitive. It also differs from GBDT
training systems by treating tree-derived path regions as candidate features and
reducing their scores on device without materializing dense membership.

\section{Future Work}

\begin{enumerate}
  \item \textbf{Buffered first-hit statistics.} Replace per-hit OptiX-side
  floating atomics with a hit-path buffer and a separate CUDA reduction, then
  compare performance and numerical behavior. This may reduce any-hit pressure
  but adds memory traffic.
  \item \textbf{Additional metrics.} Extend compact RT decision-path scoring to
  MI and Spearman only after backend parity is proven.
  \item \textbf{Real ensemble extraction.} Evaluate paths extracted from
  trained GBDT ensembles on real tabular datasets, not only synthetic partition
  grids.
  \item \textbf{GPU architecture sweep.} Compare Turing, Ampere, Ada, Hopper,
  and Blackwell RT behavior and update shape hints per architecture.
  \item \textbf{Graph replay.} Capture repeated resident RT scoring calls where
  geometry and buffer shapes remain stable.
  \item \textbf{Hardware profiling.} Automate NCU collection for visible SM,
  memory, branch, launch, and RT-adjacent counters.
\end{enumerate}

\section{Conclusion}

This report shows that a subset of GBDT-derived decision-path scoring can be
mapped effectively onto hardware ray tracing cores. The strongest case is
finite bounded non-overlapping 2D path partitions: the same structure that makes
a decision tree assign each sample to one leaf also makes first-hit RT traversal
semantically valid. By keeping the public result compact and reducing scores on
device, the implementation avoids the memory cost of materialized membership.

The current GAFIME CUDA implementation reaches approximately 2.3--2.45 trillion
membership-equivalent evaluations/s on a parity-covered RTX 4060 Laptop
workload, and 13.7 trillion evaluations/s on a larger throughput-only workload.
The result is not a general claim that RT cores accelerate all GBDT operations.
It is a narrower systems result: when decision-path candidates are bounded,
finite, low-dimensional, non-overlapping partitions, RT traversal plus on-device
compact reduction can be an unusually effective scoring primitive.

\begin{thebibliography}{10}

\bibitem{friedman2001gbm}
Jerome H. Friedman.
\newblock Greedy function approximation: A gradient boosting machine.
\newblock \emph{The Annals of Statistics}, 29(5):1189--1232, 2001.
\newblock \url{https://doi.org/10.1214/aos/1013203451}.

\bibitem{chen2016xgboost}
Tianqi Chen and Carlos Guestrin.
\newblock XGBoost: A scalable tree boosting system.
\newblock In \emph{Proceedings of the 22nd ACM SIGKDD International Conference
on Knowledge Discovery and Data Mining}, pages 785--794, 2016.
\newblock \url{https://doi.org/10.1145/2939672.2939785}.

\bibitem{ke2017lightgbm}
Guolin Ke, Qi Meng, Thomas Finley, Taifeng Wang, Wei Chen, Weidong Ma, Qiwei Ye,
and Tie-Yan Liu.
\newblock LightGBM: A highly efficient gradient boosting decision tree.
\newblock In \emph{Advances in Neural Information Processing Systems}, 30,
2017.
\newblock \url{https://papers.nips.cc/paper/6907-lightgbm-a-highly-efficient-gradient-boosting-decision-tree}.

\bibitem{prokhorenkova2018catboost}
Liudmila Prokhorenkova, Gleb Gusev, Aleksandr Vorobev, Anna Veronika Dorogush,
and Andrey Gulin.
\newblock CatBoost: Unbiased boosting with categorical features.
\newblock In \emph{Advances in Neural Information Processing Systems}, 31,
2018.
\newblock \url{https://papers.nips.cc/paper/7898-catboost-unbiased-boosting-with-categorical-features}.

\bibitem{parker2010optix}
Steven G. Parker, James Bigler, Andreas Dietrich, Heiko Friedrich, Jared
Hoberock, David Luebke, David McAllister, Morgan McGuire, Keith Morley, Austin
Robison, and Martin Stich.
\newblock OptiX: A general purpose ray tracing engine.
\newblock \emph{ACM Transactions on Graphics}, 29(4), 2010.
\newblock \url{https://doi.org/10.1145/1778765.1778803}.

\bibitem{nvidia2018turing}
NVIDIA.
\newblock NVIDIA Turing GPU Architecture.
\newblock White paper, 2018.
\newblock \url{https://www.nvidia.com/content/dam/en-zz/Solutions/design-visualization/technologies/turing-architecture/NVIDIA-Turing-Architecture-Whitepaper.pdf}.

\bibitem{aila2009raytraversal}
Timo Aila and Samuli Laine.
\newblock Understanding the efficiency of ray traversal on GPUs.
\newblock In \emph{Proceedings of High Performance Graphics}, 2009.
\newblock \url{https://users.aalto.fi/~ailat1/HPG2009/}.

\bibitem{aila2012raytraversal}
Timo Aila, Samuli Laine, and Tero Karras.
\newblock Understanding the efficiency of ray traversal on GPUs: Kepler and
Fermi addendum.
\newblock NVIDIA Technical Report NVR-2012-02, 2012.
\newblock \url{https://research.nvidia.com/publication/2012-06_understanding-efficiency-ray-traversal-gpus-kepler-and-fermi-addendum}.

\bibitem{wald2014embree}
Ingo Wald, Sven Woop, Carsten Benthin, Gregory S. Johnson, and Manfred Ernst.
\newblock Embree: A kernel framework for efficient CPU ray tracing.
\newblock \emph{ACM Transactions on Graphics}, 33(4), 2014.
\newblock \url{https://doi.org/10.1145/2601097.2601199}.

\end{thebibliography}

\end{document}
